\scanpage{65}
\uline{formule de Jacobi} Soient $a,b,\mathord{\text{---}}\ l$ les poles \deleted{d'une} \inserted{d'une} fonction elliptique de \unsure{pér|pés} $2\omega_1$ et $2\omega_2$ (distincts à \unsure{$2\omega$|$1\omega$} près) $a_1\ b_1\ \mathord{\text{---}}\ l_1$ les zéros ; on a
\[
a+b\mathord{---}+l=a_1\unsure{+|-}b_1\mathord{---}+l_1+\Omega\ .\quad \Omega\ \text{est une}
\]

On a alors
\[
f(z)=c^{\mathrm{te}}\times\frac{\sigma(z-a_1)\sigma(\unsure{z|B}-b_1)\mathord{--}\sigma(z-l_1-\Omega)}
{\sigma(z-a)\sigma(z-b)\mathord{---}\sigma(z-l)}
\]
soit
\[
f(z)=c^{\mathrm{te}}\times\frac{\sigma(z-a_1)\sigma(z-b_1)\mathord{-}\sigma(z-l_1)}
{\sigma(z-a)\sigma(z-b)\mathord{---}\sigma(z-l+\Omega)}
\]
soit \unsure{aussi|enfin|encore}, enfin
\[
f(z)=c^{\mathrm{te}}\times e^{2(n_1\eta_1+n_2\eta_2)z}
\frac{\sigma(z-a_1)\mathord{--}\sigma(z-l_1)}{\sigma(z-a)\mathord{--}\sigma(z-l)}
\]
en posant $\Omega=2n_1\omega_1+2n_2\omega_2$

\uline{formule d'addition}

\[
\begin{aligned}
\unsure{p|\wp}u-p\unsure{v|w}&=-\frac{\sigma(u+v)\sigma(u-v)}{\sigma^2u\,\sigma^2v}\\
\zeta(u+v)&=\zeta u+\zeta v+\frac12\,\frac{p'u-p'v}{pu-pv}\\
p(u+\unsure{u|v})+pu+pv&=\frac14\left(\frac{p'u-p'v}{pu-pv}\right)^2=[\zeta(u+v)-\zeta u-\zeta v]^2
\end{aligned}
\]
% archive sheet number: 64
